% 阴影面积的意义——物理量的累积效应
% keys 速度、电荷量、平均物理量
% license Usr
% type Tutor



\subsection{物理量的累积效应}

\subsubsection{速度}
在初中里，我们已经学过“速度”，也就是单位时间内的路程。然而这个概念的局限很多，一方面，路程和时间都是数字，“速度”很自然的，也没有给出其他更多的信息。另一方面，当我们用初中定义去描述任何对象的快慢时，都是默认其快慢不变的。既然小明从西到东的速度是$6 \mathrm{m/s}$，那么在这个过程中，在任意的$0.02\mathrm s$内，移动的路程便总是$1.2\mathrm m$，这显然是不够贴近实际的。用“位移”代替“路程”，我们便可以解决第一个问题，这样定义的速度我们称之为平均速度。新增了方向信息后，那如何描述任意时刻的快慢呢？

回想一下，在我们用“路程-时间”图像描述物体运动时，匀速直线运动是一条斜率不变的直线，“匀速”便是路程均匀变化。受此启发，我们不妨用\textbf{斜率}来表示物体在某个时刻的速度，称之“\textbf{瞬时速度}”。这也是因为对于位移-时间图像，你总可以作出函数图像在任意一点的切线，该切线总是唯一的，所以定义切线斜率就是此时此刻的瞬时速度，显然，瞬时速度是平均速度在某个时刻附近的近似。
%瞬时速度和平均速度
当速度均匀变化时，我们称之为匀变速直线运动，并把速度-时间图像里的斜率称之为加速度。物体并非总是做匀变速直线运动。更符合实际的情况是变加速，即加速度总是变化的运动。
%匀变速和变加速的时间图像

匀速直线运动的速度时间图像是一条平行于时间轴的直线。对于任意一个时间段，时间区间与速度的乘积显然就是位移。但对于一般的速度-时间图像，这就很难成立了。从定义出发，你不知道哪个是特殊的“平均速度”，毕竟，平均速度其实是知晓位移后导出的，你怎可本末倒置地去找这个特殊值呢？但我们可以据匀速的情况猜测，速度函数与两个时间轴之间的阴影面积就是这段时间的的位移。

为了得到阴影面积，我们可以用矩形集合来近似。将时间段划分为区间长度一致的小时间段，并取对应时间段内的任意一个函数值作为矩形的高。如下图所示，随着区间长度逼近0，矩形的面积之和也就逼近阴影面积，又因为当时间足够小时，瞬时速度近似于平均速度，因而对应矩形面积近似于该段位移，则时间取得足够小时，矩形面积之和，即总位移约等于阴影面积的逼近值。当我们说“速度-时间图像的阴影面积是位移”时，实际上说的是逼近值。


\begin{figure}[ht]
\centering
\includegraphics[width=14cm]{./figures/9254f1a64edc702f.png}
\caption{用矩形集合近似阴影面积} \label{fig_area_5}
\end{figure}

速度是位移-时间图像的斜率，速度-时间图像的阴影面积是位移，这是巧合么？我们能否推而广之，认为若$F(x)$的斜率是$y(x)$，则$y(x)$在$x_1,x_2$之间的阴影面积便是$\Delta F=F(x_2)-F(x_1)$。本科阶段的微积分基本定理可以证明这一点，此处暂不赘述。
\subsubsection{其他运动学量}
\begin{itemize}
\item 加速度-时间图像 

加速度是速度函数的斜率，因此加速度对时间的阴影面积，或言对时间的积累，便是这段时间的速度变化量。
\footnote{因为矢量遵循平行四边形定则，因此即使本节推导是一维下进行的，也能自然推广到一个多维运动。}
\item 力-位移图像 

当合外力不变时，$F\Delta x=W$，任意力对位移的积累是该力在这段位移中所做的功。
\item 合力-时间图像

由牛顿定律可知，$F\Delta t=ma\Delta t=m\Delta v=\Delta(mv)$，定义$mv$为动量，则合力对时间的积累便是动量变化量，这就是动量定理。
\end{itemize}

\subsubsection{电磁学}
\begin{itemize}
\item 电流-时间图像 因为电流是电荷量-时间图像的斜率，因此电流对时间的积累便是这段时间通过对应截面的电荷量。
\item 电场-位移图像（粒子只受电场力）
因为$qE(x_B-x_A)=q(\phi_A-\phi_B)$，所以电场对位移的积累是电势差，读者需要注意电势差的正负号。
\item 磁感应强度-面积图像

同理，由于$B\Delta S=\Delta \phi$，所以磁感应强度对线圈面积的积累是该线圈在对应过程通过的磁通量。
\end{itemize}

\subsubsection{热学}
压强-体积图像，对于中学阶段常见的气体模型，因为$p\Delta V=W$，即系统气体对容器边界所做的功，所以压强对体积的积累便是这部分功。
\subsection{“平均物理量”}
在高中物理学习期间，我们接触的第一个平均物理量是平均速度。平均速度采取了一个无论物体做什么运动都适用的定义，即$\bar v\Delta t=x$。对于任意运动，我们只要求出了速度-时间图像的阴影面积，即位移，就能自然求出平均速度。对于一个匀变速直线运动，我们可以利用位移公式，证明任意一段时间内的平均速度便是首末速度$(v_1,v_2)$的算术平均值，即$(v_1+v_2)/2$。

实际上，用阴影面积除以自变量变化值也是目前最广泛的平均物理量的定义，如平均电流，冲量中的平均力，等等。在此定义下，首末物理量平均值作为平均物理量仅适用于\textbf{线性变化}的物理量。这甚至不需要证明。回想一下，我们求平均速度的时候，依据的位移面积是根据线性变化的速度求得的。自然而然，把速度-时间替换成任意一个物理量和自变量，求得的面积也只是把初速度，加速度，时间替换成初物理量，物理量对自变量的斜率，自变量。读者可自行验证，弹簧弹力的平均值亦是首末弹力的平均值。
